\scanpage{53}
\margin{\uline{propriétés} \uline{non topiques} \inserted{(pour topos)} absolues pour anneaux $\left\{\begin{array}{l}\text{artinien, noethérien, de t.f. sur $\mathbb Z$ (ou sur}\\ \text{$n$ composantes connexes\qquad $\unsure{\mathbb Q|Q}$, \deleted{\unsure{ou sur}}}\\ \text{$n$ composantes irréductibles\qquad ou $\mathbb Z S^{-1}\ldots$)}\end{array}\right.$}

\deleted{\uline{Propriétés relatives d'anneaux}}

Anneaux plongeables dans des corps : il \uline{faut} qu'ils soient intègres,
\margin{(voir anneau \unsure{hensélien}, ou ---- \unsure{hens\ldots})}
mais ce n'est pas suffisant \inserted{même s'il est local \unsure{hensélien|hypothèse}}. Pour l'anneau intègre universel, les sections locales qui sont partout non nulles (donc qui sont $A$-régulières) sont celles qui sont inversibles : l'anneau total des fractions \deleted{\unsure{doit}} est égal à $A$ lui-même. Une condition nécessaire pour qu' \deleted{\unsure{A soit un}} un Anneau se plonge \unsure{$\to X$|$X'\to X$} dans un corps, c'est que la section nulle (donc tte autre section) soit \unsure{sommande|sommante} direct de $A$, et c'est sans doute suffisant

\uline{Modules sur anneau}
\[
 X_{\mathrm{mod}}=\widehat{\mathit{Mod}^{\circ}_{pf}}
\]
\[
 \begin{aligned}\mathit{Mod}_{pf}&=\text{catégorie des modules de prés. finie sur $\mathbb Z$-algèbres de prés finie $\unsure{\in|\mathscr E}$}\\
 \mathit{Mod}_{pf}^{\circ}&\simeq\text{\deleted{schémas} Modules \deleted{\unsure{fibrés}} sur schémas affines de t.f. sur $\mathbb Z$.}
 \end{aligned}
\]
\diagram{\begin{tikzcd}[ampersand replacement=\&,row sep=small]
 M\arrow[r]\arrow[d,no head]\&N\arrow[d,no head]\\
 A\arrow[r]\&B\\
 \mathscr F\&\mathscr G\\
 Y\& X\arrow[l]
\end{tikzcd}}
\diagram{$\begin{array}{l}
 \left\{\begin{array}{l}\text{Localement libre}\\ \text{libre, projectif / injectif (?)}\end{array}\right.\\
 \left\{\begin{array}{l}\text{localemt libre de rang $n$}\\\text{libre de rang $n$}\end{array}\right.\quad(n\geq 1)\\
 \text{\unsure{ou|$n$}}\\ \text{de \unsure{t.f.|N.f.}}\\ \text{de prés finie}\\ \text{monogène}
 \end{array}$}
\margin{ne passent pas à la limite inductive ? donc pas topiques}% slanted

\deleted{\unsure{Plats, libres, projectifs}}
\[
 M\ \text{inv.}\quad\Longleftrightarrow\quad M\otimes_A\underbrace{\underline{\operatorname{Hom}}_A(M,A)}\xrightarrow{\sim}A
\]
\margin{n'est pas exprimable par \unsure{lim|les lim|$\varinjlim$} / ne commute pas au chgt de base du \unsure{sg}}% slanted

\deleted{\unsure{$\exists$|$\exists N$} $M\otimes N\simeq A$ \quad $X\longmapsto M\otimes X$\quad $\unsure{\operatorname{Hom}(X,Y)\simeq\operatorname{Hom}(TX,TY)|\underline{\operatorname{Hom}}(X,Y)\to\underline{\operatorname{Hom}}(f^*X,f^*Y)}$\quad $\unsure{\operatorname{Hom}(T'TX,T'Y)}$}

\deleted{\unsure{$F(M\otimes$}$\,\unsure{N)}$ \quad \unsure{$T_M T_N\simeq 1\Longrightarrow$} \unsure{$\mathcal U$} \quad $T_M T_N\simeq\mathrm{id}$\quad donc $T_N,T_M$ des équivalences\quad $\unsure{\operatorname{Hom}(X,Y)|\underline{\operatorname{Hom}}(X,Y)}$}

\deleted{\unsure{mor} $M'\to M$ et $M\otimes P\to M\otimes Q$\quad $T_N T_M\simeq\mathrm{id}$\quad donc exacts,}

\deleted{donc $M$ et $N$ \inserted{fid} plats. $M=\lim M_i$, les $M_i\subset M$ de t.f., donc $A\simeq M\otimes N=\underset{\longrightarrow}{\lim}\,M_i\otimes N$, donc $M_i\otimes N\subset A$, donc $\unsure{\exists}$ on a $M_i\otimes N=A$ pour \unsure{$i$|$A'$|$A$} grand, donc $M_i=M$, donc $M$ de t.f., \deleted{\unsure{Il est}} $N$ pareil. $M$ est monogène}
% archive sheet number: 52
